\begin{figure}[htbp]
\centering
\begin{tikzpicture}[font=\small]
  \fill[blue!7] (-1.25,0) ellipse (2.05 and 1.45);
  \fill[orange!9] (1.25,0) ellipse (2.05 and 1.45);
  \draw[blue!60!black,thick] (-1.25,0) ellipse (2.05 and 1.45);
  \draw[orange!70!black,thick] (1.25,0) ellipse (2.05 and 1.45);
  \node[blue!60!black] at (-2.4,1.35) {$U_e$};
  \node[orange!70!black] at (2.4,1.35) {$U_f$};
  \node at (-1.7,.75) {$(w,j)$};
  \node at (-2.0,-.1) {$(v,l)$};
  \node at (-1.45,-.85) {$(w,l)$};
  \node at (1.7,.75) {$(z,j)$};
  \node at (2.0,-.1) {$(v,m)$};
  \node at (1.45,-.85) {$(z,m)$};
  \fill[linkblue] (0,.55) circle (1.7pt);
  \node[align=center] at (0,-.05) {$\kappa$\\$=(v,j)$};
  \node at (0,-1.9) {$M_\kappa=U_e\cup U_f,\qquad
                    \Gamma_\kappa=M_\kappa\setminus\{\kappa\}$};
\end{tikzpicture}
\caption{The seven coordinates available at the corner
$\kappa=(v,j)$, where $j=\{v,w,z\}$, $e=\{v,w\}$ and
$f=\{v,z\}$. The opposite faces on $e$ and $f$ are $l$ and $m$.
The two activation branches keep the same six outer coordinates;
one also keeps the common corner marked by the dot.}
\label{fig:kernel-corner-activation}
\end{figure}
